% Optimization chapter skeleton. Requires \TemplatePrompt from the selected paper family.
\section{优化模型的建立与求解}

\subsection{任务目标与模型选择}
\TemplatePrompt{定义决策变量、目标、资源约束和决策时域；说明为何选择该优化模型，并给出被排除的替代方案。}

\subsection{可运行基线}
\TemplatePrompt{给出与主模型同输出类型的可行基线，例如当前策略、规则策略或松弛模型，并统一目标函数和约束检查。}

\subsection{数学模型}
\TemplatePrompt{逐项给出目标函数、硬约束、软约束、变量域、参数来源和单位；说明线性化或近似的影响。}

\subsection{算法与停止条件}
\TemplatePrompt{记录求解器、版本、容差、时间上限、随机种子和不可行问题的诊断方式。}

\subsection{主结果与可行性证据}
\TemplatePrompt{比较主模型与基线的目标值、约束余量、运行时间和重复运行波动；图表只读取实验产物。}

\subsection{灵敏度、稳健性与结论边界}
\TemplatePrompt{扰动关键成本、容量和需求参数；报告策略是否改变、可行率和性能代价，并限定外推范围。}
